\section{Resumen}

Esta tesis examina si los grandes modelos de lenguaje (LLMs) satisfacen 
los criterios requeridos para la inferencia causal, o si únicamente producen asociaciones estadísticas. 
La hipótesis es que los LLMs no satisfacen los criterios requeridos para la inferencia causal. 
Debido a que operan sobre una distribución observacional comprimida en \textit{tokens} que les 
permite producir respuestas casualmente plausibles en el nivel de 
asociación (L1), pero carecen de la representación mecanística que 
exigen los niveles de intervención (L2) y contrafactual (L3) de la 
jerarquía de \textcite{pearl2009causality}. El trabajo integra dos componentes.

El \textbf{primero} es un marco teórico que articula, a partir de la relación de 
modelado de \textcite{rosen2012anticipatory}, el modelo causal de 
\textcite{pearl2009causality} y el principio de mecanismos causales 
independientes de \textcite{scholkopf2021toward}, un criterio formal para 
distinguir inferencia causal de simulación predictiva. El argumento 
central es que la distinción no se resuelve con una optimización sino que viene 
desde la arquitectura; un sistema con comprensión estructural dispone de una representación  
\textit{mecanística}, la capacidad de un sistema de representar internamente 
las funciones estructurales que determinan el valor de cada variable 
a partir de sus causas,  $X_i := f_i(\text{PA}_i, U_i)$ que componen el modelo causal estructural (SCM) \parencite{pearl2000causality}. 
En cambio, un LLM dispone de una distribución de probabilidad sobre tokens optimizada para predecir texto 
humano.

El \textbf{segundo} es un componente computacional que implementa un programa de 
evaluación sobre el benchmark CLadder \parencite{jin2023cladder}, diseñado 
explícitamente sobre la jerarquía causal de \textcite{pearl2009causality} 
para evaluar los tres niveles de razonamiento causal en modelos de lenguaje: asociación (L1), 
intervención (L2) y contrafactual (L3). El programa consulta LLMs mediante 
API, clasifica sus respuestas según el nivel causal requerido, y construye 
un grafo acíclico dirigido (DAG) de patrones de fallo que visualiza en qué 
niveles de la jerarquía causal fallan.  
La implementación utiliza \texttt{networkx} para la construcción del grafo 
y la API de un LLM  para las consultas. 

Los resultados empíricos se interpretaran con apoyo del marco teórico para 
argumentar que los fallos observados no son aleatorios sino estructuralmente 
consistentes con la hipótesis principal: los LLMs se concentran 
sistemáticamente en el nivel de asociación (L1) de la jerarquía de Pearl y exhiben 
un fallo en las relaciones de intervención (L2) y contrafactual (L3). 